\documentclass[10pt, a4paper]{article}

\usepackage[T1]{fontenc}
\usepackage[utf8]{inputenc}
\usepackage[english]{babel}
\usepackage{charter}
\usepackage[top=1.4cm, bottom=1.4cm, left=0.9cm, right=1.4cm]{geometry}
\usepackage{xcolor}
\usepackage{enumitem}
\usepackage{hyperref}
\usepackage{array}
\usepackage{tabularx}
\usepackage{microtype}
\usepackage{fontawesome5}
\usepackage{paracol}
\usepackage{eso-pic}
\usepackage{needspace}
\usepackage{ragged2e}
\setlength{\JustifyingParindent}{0pt}
\setlength{\RaggedRightParindent}{0pt}

%% COLORES
\definecolor{accent}{HTML}{1F4E79}
\definecolor{lightaccent}{HTML}{2E75B6}
\definecolor{gray}{HTML}{555555}
\definecolor{lightgray}{HTML}{888888}
\definecolor{rulegray}{HTML}{CCCCCC}
\definecolor{sidebg}{HTML}{EEF3F8}
\definecolor{tagbg}{HTML}{DCE6F1}

\hypersetup{colorlinks=true, urlcolor=lightaccent, linkcolor=lightaccent,
  pdftitle={Lara Mar\'ia V\'azquez Gonz\'alez - CV}, pdfauthor={Lara Mar\'ia V\'azquez Gonz\'alez}}

\pagestyle{empty}
\setlength{\parindent}{0pt}
\setlength{\parskip}{0pt}

%% BARRA LATERAL: ancho de la franja de color y de la columna izquierda
\newlength{\sidebarwidth}
\setlength{\sidebarwidth}{6.0cm}
\AddToShipoutPictureBG{%
  \AtPageLowerLeft{\color{sidebg}\rule{\sidebarwidth}{\paperheight}}}
\setcolumnwidth{4.75cm, \dimexpr\textwidth-4.75cm-1.25cm\relax}
\setlength{\columnsep}{1.25cm}

\setlist[itemize]{leftmargin=1.1em, itemsep=2pt, topsep=2pt, parsep=0pt}

%% MACROS
% Titulo de seccion de la columna principal
\newcommand{\mainsection}[1]{%
  \vspace{9pt}{\large\bfseries\color{accent}#1}\par
  \vspace{2pt}{\color{rulegray}\hrule}\vspace{5pt}}
% Titulo de seccion de la barra lateral
\newcommand{\sidesection}[1]{%
  \ifhmode\unskip\unskip\fi\par\needspace{5\baselineskip}\vspace{11pt}{\normalsize\bfseries\color{accent}\MakeUppercase{#1}}\par
  \vspace{5pt}{\color{accent}\hrule width 1.2cm height 1pt}\vspace{6pt}}
% Linea de contacto con icono
\newcommand{\contact}[2]{%
  \makebox[1.2em][c]{\color{accent}#1}\hspace{2pt}{\footnotesize #2}\par\vspace{3pt}}
% Etiqueta de habilidad
\newcommand{\tag}[1]{%
  {\setlength{\fboxsep}{2pt}\colorbox{tagbg}{\footnotesize\color{accent!90!black}\vphantom{Ag}#1}}\hspace{1pt plus 2pt}}
% Entrada de experiencia: puesto, empresa, fechas, lugar
\newcommand{\cventry}[4]{%
  \begin{tabularx}{\linewidth}{@{}Xr@{}}
    \textbf{#1} & \textcolor{lightaccent}{\small #3} \\
    \textcolor{gray}{#2} & \textcolor{lightgray}{\small #4}
  \end{tabularx}\par\vspace{1pt}}
% Entrada de formacion en la barra lateral
\newcommand{\eduentry}[3]{%
  {\small\bfseries #1}\par
  {\footnotesize\color{gray}#2}\par
  {\footnotesize\color{lightaccent}#3}\par\vspace{4pt}}
% Publicacion
\newcommand{\cvpub}[5]{%
  \item \textcolor{gray}{\small #1 (#2).} \textbf{\small #3.} \textit{\small #4.} \textcolor{lightaccent}{\small #5}}

\begin{document}
\begin{paracol}{2}

%% ===================== BARRA LATERAL =====================
\raggedright\hyphenpenalty=10000

\vspace*{2pt}
\sidesection{Contact}
\contact{\faEnvelope}{\href{mailto:vazquezgonzalezlara@gmail.com}{vazquezgonzalezlara@gmail.com}}
\contact{\faMapMarker*}{Santiago de Compostela}
\contact{\faGlobe}{\href{https://laravg.github.io}{laravg.github.io}}

\sidesection{Languages}
{\small\textbf{English}\hfill C2 (Cambridge)}\par\vspace{3pt}
{\small\textbf{Spanish}\hfill Native}\par\vspace{3pt}
{\small\textbf{Galician}\hfill Native}\par

\sidesection{Education}
\eduentry{PhD in Information Technology Research}{Universidade de Santiago de Compostela}{2020 -- 2025}
\eduentry{MSc in Bioinformatics and Biostatistics}{Universitat Oberta de Catalunya}{2019 -- 2020}
\eduentry{BSc in Computer Engineering}{Universidade de Santiago de Compostela}{2015 -- 2019}

\sidesection{Programming}
\tag{Python} \tag{R} \tag{Bash} \tag{SQL}

\sidesection{Sequencing}
\tag{QIIME 2} \tag{DADA2} \tag{BLAST} \tag{Illumina}

\sidesection{R and statistics}
\tag{phyloseq} \tag{vegan} \tag{ggplot2} \tag{R Markdown} \tag{mixed models} \tag{differential abundance}

\sidesection{Python and ML}
\tag{pandas} \tag{NumPy} \tag{Biopython} \tag{scikit-learn} \tag{matplotlib} \tag{ensemble learning} \tag{feature selection}

\sidesection{Environments}
\tag{Linux} \tag{HPC/SLURM} \tag{conda} \tag{Docker} \tag{Git} \tag{GitHub} \tag{GitLab}

\sidesection{Currently learning}
{\small\bfseries Nextflow}\par
{\footnotesize\color{gray}Courses: Hello Nextflow and Nextflow Run (Seqera, 2026)}\par

\sidesection{Profiles}
\contact{\faLinkedin}{\href{https://linkedin.com/in/laramariavazquezgonzalez}{laramariavazquezgonzalez}}
\contact{\faGithub}{\href{https://github.com/laravg}{github.com/laravg}}
\contact{\faGitlab}{\href{https://gitlab.citius.gal/users/lara.vazquez/projects}{gitlab.citius.gal/lara.vazquez}}
\contact{\faOrcid}{\href{https://orcid.org/0000-0002-0122-5266}{0000-0002-0122-5266}}
\contact{\faGraduationCap}{\href{https://scholar.google.com/citations?user=F9JtFRwAAAAJ}{Google Scholar}}

%% ===================== COLUMNA PRINCIPAL =====================
\switchcolumn
\justifying

{\fontsize{24}{28}\selectfont\bfseries\color{accent} Lara Mar\'ia V\'azquez Gonz\'alez}\par
\vspace{5pt}
{\large\color{gray} Bioinformatics \textperiodcentered{} NGS Data Analysis \textperiodcentered{} Machine Learning}\par
\vspace{6pt}
{\color{accent}\rule{\linewidth}{1.5pt}}\par

\mainsection{Profile}
Bioinformatician with a PhD in Information Technology Research (USC, 2025), an MSc in Bioinformatics and Biostatistics and a degree in Computer Engineering. I work on NGS data analysis, mainly 16S rRNA amplicon sequencing and clinical metagenomics, building end-to-end analyses (quality control, preprocessing, statistical analysis and reporting) in Python, R and Bash on Linux and HPC environments with SLURM. I have developed open-source bioinformatics tools and applied machine learning to microbiome data. Six scientific publications, three as first author in Q1 journals.

\mainsection{Experience}

\cventry{Collaborating Researcher}{Universidade de Santiago de Compostela}{Sep 2026 -- Present}{Santiago de Compostela}
\begin{itemize}
  \item Processing of 16S rRNA amplicon sequencing data (V3-V4, Illumina paired-end) from clinical oral microbiome studies, combining several sequencing runs: quality control, primer trimming (cutadapt), per-run DADA2 denoising, taxonomic assignment (VSEARCH), contaminant removal with negative controls (decontam) and phylogeny (MAFFT, FastTree) in QIIME 2, with version-controlled scripts.
  \item Statistical analysis in R (phyloseq, vegan, ANCOM-BC2, glmmTMB): alpha and beta diversity, PERMANOVA, differential abundance and mixed models for paired designs.
  \item Methods and results reports (R Markdown and LaTeX) and documented deliverables for clinical research teams.
\end{itemize}
\vspace{6pt}

\cventry{AI Research Technician}{CESGA, Galicia Supercomputing Center}{Jun 2025 -- Jun 2026}{Santiago de Compostela}
\begin{itemize}
  \item Developed AIRA16S during a research stay at CiTIUS (USC): a Python pipeline for the automated inventory of human 16S rRNA sequencing projects in public databases (NCBI SRA, BioProject, PubMed, Europe PMC), with LLM-based metadata extraction run on an HPC cluster, semantic validation and an interactive results viewer.
  \item OneHealth DataSpace project (NextGenerationEU), following Scrum: first version of the project website (Next.js, TypeScript, Git) and support for the deployment of the data space platform.
\end{itemize}
\vspace{6pt}

\cventry{Predoctoral Researcher}{USC, CiTIUS \textperiodcentered{} Xunta de Galicia grant ED481A-2021}{Oct 2020 -- May 2025}{Santiago de Compostela}
\begin{itemize}
  \item PhD thesis (BIOMETA) on bioinformatics procedures for microbiome analysis through 16S rRNA sequencing, applied to the oral microbiome, with model training and execution on an HPC cluster via SLURM.
  \item Developed three open-source tools, version-controlled with Git and published on GitLab, each with a first-author Q1 publication: PrimerEvalPy, 16SGOSeq and EPheClass.
  \item International research stay at HEIG-VD, Switzerland (2022), where I developed the first version of EPheClass.
\end{itemize}
\vspace{6pt}

\cventry{Software Developer}{Altia}{Mar 2019 -- Sep 2020}{Santiago de Compostela}
{\small\color{gray} Software development in an industry environment, first as an intern.}\par

\mainsection{Software}
\begin{itemize}[itemsep=3pt]
  \item \textbf{PrimerEvalPy}: \textit{in silico} evaluation of primers and primer pairs against sequence databases, with coverage metrics at each taxonomic level. \href{https://gitlab.citius.gal/lara.vazquez/PrimerEvalPy}{\small [GitLab]}
  \item \textbf{16SGOSeq}: curated database of 16S rRNA genes and their variants in complete genomes of oral bacteria and archaea, with taxonomy from domain to strain. \href{https://gitlab.citius.gal/lara.vazquez/16sgoseq}{\small [GitLab]}
  \item \textbf{EPheClass}: phenotype classification from 16S rRNA count tables with a dynamic ensemble of machine learning models (kNN, Random Forest, SVM, XGBoost, MLP); ROC AUC 0.974. \href{https://gitlab.citius.gal/lara.vazquez/epheclass}{\small [GitLab]}
  \item \textbf{Oral16SAmpliconSimilarity}: \textit{in silico} analysis of 16S rRNA amplicon similarity across oral prokaryotic genomes, to identify which species become indistinguishable depending on the primer pair used. \href{https://gitlab.citius.gal/lara.vazquez/oral16sampliconsimilarity}{\small [GitLab]}
  \item \textbf{AIRA16S}: automated inventory of human 16S rRNA sequencing projects in public databases, with LLM-based metadata extraction and semantic validation, integrated into a health data space.
\end{itemize}

\mainsection{Publications}
{\sloppy
\begin{itemize}[itemsep=3pt]
  \cvpub{V\'azquez-Gonz\'alez, L. et al.}{2024}{PrimerEvalPy: a tool for in-silico evaluation of primers for targeting the microbiome}{BMC Bioinformatics (Q1)}{\href{https://doi.org/10.1186/s12859-024-05805-7}{doi:10.1186/s12859-024-05805-7}}
  \cvpub{V\'azquez-Gonz\'alez, L. et al.}{2025}{A curated bacterial and archaeal 16S rRNA Gene Oral Sequences dataset}{Scientific Data (Q1)}{\href{https://doi.org/10.1038/s41597-025-05050-4}{doi:10.1038/s41597-025-05050-4}}
  \cvpub{V\'azquez-Gonz\'alez, L. et al.}{2025}{EPheClass: ensemble-based phenotype classifier from 16S rRNA gene sequences}{Frontiers in Bioinformatics (Q1)}{\href{https://doi.org/10.3389/fbinf.2025.1514880}{doi:10.3389/fbinf.2025.1514880}}
  \cvpub{Regueira-Iglesias, A., V\'azquez-Gonz\'alez, L. et al.}{2023}{In silico evaluation and selection of the best 16S rRNA gene primers for use in next-generation sequencing to detect oral bacteria and archaea}{Microbiome (Q1)}{\href{https://doi.org/10.1186/s40168-023-01481-6}{doi:10.1186/s40168-023-01481-6}}
  \cvpub{Regueira-Iglesias, A., V\'azquez-Gonz\'alez, L. et al.}{2023}{Impact of 16S rRNA Gene Redundancy and Primer Pair Selection on the Quantification and Classification of Oral Microbiota in Next-Generation Sequencing}{Microbiology Spectrum (Q2)}{\href{https://doi.org/10.1128/spectrum.04398-22}{doi:10.1128/spectrum.04398-22}}
  \cvpub{Regueira-Iglesias, A., V\'azquez-Gonz\'alez, L. et al.}{2022}{In-Silico Detection of Oral Prokaryotic Species With Highly Similar 16S rRNA Sequence Segments Using Different Primer Pairs}{Frontiers in Cellular and Infection Microbiology (Q1)}{\href{https://doi.org/10.3389/fcimb.2021.770668}{doi:10.3389/fcimb.2021.770668}}
\end{itemize}
}

\mainsection{Conferences}
{\sloppy
\begin{itemize}[itemsep=3pt]
  \cvpub{V\'azquez-Gonz\'alez, L. et al.}{2023}{An Ensemble-Based Phenotype Classifier to Diagnose Crohn's Disease from 16S rRNA Gene Sequences}{11th Iberian Conference on Pattern Recognition and Image Analysis (IbPRIA), Alicante. Poster}{\href{https://doi.org/10.1007/978-3-031-36616-1_44}{doi:10.1007/978-3-031-36616-1\_44}}
  \cvpub{V\'azquez-Gonz\'alez, L. et al.}{2021}{PrimerEvalPy: a tool for in-silico evaluation of primers for targeting the microbiome}{43rd Annual International Conference of the IEEE Engineering in Medicine \& Biology Society (EMBC), online. One-page paper}{}
\end{itemize}
}

\end{paracol}
\end{document}
